% GENERATED FILE. Edit canonical lessons, then run npm run generate:books.
% canonical-source-hash: fnv1a64:850826de74e1d6d6
% canonical-lessons: IT-C491-cotta, IT-C491-corteggiare, IT-C491-innamorato, IT-C491-tenerezza, IT-C491-bacio

\chapter{Innamorati}
\label{ch:it-innamorati}
\hlchaptermodality{491}

\begin{chapteropening}
\textbf{By the end of this chapter:} \emph{I can talk about crushes, courting and falling in love.}

Complete the last lesson of chapter 491: I can talk about crushes, courting and falling in love.
\end{chapteropening}

\section[la cotta]{la cotta — crush (on someone)}

\label{lesson:IT-C491-cotta}



\begin{quote}
Before the new one: say the Italian for to yawn, then the Italian for to stretch.
\end{quote}

\subsection*{You'll want to know: la cotta}
\textbf{la cotta} — \textquotedblleft{}crush (on someone)\textquotedblright{}.

\begin{quote}
A quindici anni avevo una cotta per un compagno di classe.
\end{quote}

\begin{quote}
\emph{At fifteen I had a crush on a classmate.}
\end{quote}

\emph{Prendersi una cotta per qualcuno} — to get a crush on someone. Informal.

\begin{cousinweb}
\emph{cotta} is the feminine past participle of \emph{cuocere} \textquotedblleft{}to cook\textquotedblright{}, from Latin \emph{coquere} (Lewis \& Short): someone \emph{cotto} is cooked through with love.
\end{cousinweb}

\begin{grammarlens}[title={What you've built}]
One more word for telling, explaining and describing at B1.
\end{grammarlens}

\subsection*{Your turn}
\begin{itemize}
\raggedright
  \item \emph{Say it:} \emph{la cotta}
  \item \emph{Say it:} \emph{la cotta}, once more
  \item \emph{Say it:} a sentence of your own with \emph{la cotta}
  \item \emph{Recall:} say the Italian for to yawn, then the Italian for to stretch, then say \emph{la cotta} again
\end{itemize}

\begin{tcolorbox}[breakable,colback=teal!4,colframe=teal!35!black,title={Before you move on}]
What does la cotta mean? (\textquotedblleft{}Crush (on someone)\textquotedblright{}.) Say it once more.
\end{tcolorbox}
\section[corteggiare]{corteggiare — to court, to woo}

\label{lesson:IT-C491-corteggiare}



\begin{quote}
Before the new one: say the Italian for to stretch, then the Italian for crush.
\end{quote}

\subsection*{You'll want to know: corteggiare}
\textbf{corteggiare} — \textquotedblleft{}to court, to woo\textquotedblright{}.

\begin{quote}
L'ha corteggiata per mesi prima del primo appuntamento.
\end{quote}

\begin{quote}
\emph{He courted her for months before the first date.}
\end{quote}

Built on \emph{corte} \textquotedblleft{}court\textquotedblright{}: at first it meant to serve at a court.

\begin{cousinweb}
\emph{corteggiare} is built on \emph{corte} \textquotedblleft{}court\textquotedblright{}, inherited from Latin \emph{cohors, cohortis} \textquotedblleft{}an enclosed yard\textquotedblright{} (Lewis \& Short).
\end{cousinweb}

\begin{grammarlens}[title={What you've built}]
One more word for telling, explaining and describing at B1.
\end{grammarlens}

\subsection*{Your turn}
\begin{itemize}
\raggedright
  \item \emph{Say it:} \emph{corteggiare}
  \item \emph{Say it:} \emph{corteggiare}, once more
  \item \emph{Say it:} a sentence of your own with \emph{corteggiare}
  \item \emph{Recall:} say the Italian for to stretch, then the Italian for crush, then say \emph{corteggiare} again
\end{itemize}

\begin{tcolorbox}[breakable,colback=teal!4,colframe=teal!35!black,title={Before you move on}]
What does corteggiare mean? (\textquotedblleft{}To court, to woo\textquotedblright{}.) Say it once more.
\end{tcolorbox}
\section[l'innamorato]{l'innamorato — sweetheart; in love}

\label{lesson:IT-C491-innamorato}



\begin{quote}
Before the new one: say the Italian for crush, then the Italian for to court.
\end{quote}

\subsection*{You'll want to know: l'innamorato}
\textbf{l'innamorato} — \textquotedblleft{}sweetheart; in love\textquotedblright{}.

\begin{quote}
Gli innamorati passeggiano sul lungomare.
\end{quote}

\begin{quote}
\emph{The sweethearts stroll along the seafront.}
\end{quote}

As an adjective: \emph{è innamorato di lei} — he is in love with her.

\begin{cousinweb}
\emph{innamorato} is the past participle of \emph{innamorare} \textquotedblleft{}to make fall in love\textquotedblright{} (EtymDB), built on \emph{amore} \textquotedblleft{}love\textquotedblright{}, inherited from Latin \emph{amor}.
\end{cousinweb}

\begin{grammarlens}[title={What you've built}]
One more word for telling, explaining and describing at B1.
\end{grammarlens}

\subsection*{Your turn}
\begin{itemize}
\raggedright
  \item \emph{Say it:} \emph{l'innamorato}
  \item \emph{Say it:} \emph{l'innamorato}, once more
  \item \emph{Say it:} a sentence of your own with \emph{l'innamorato}
  \item \emph{Recall:} say the Italian for crush, then the Italian for to court, then say \emph{l'innamorato} again
\end{itemize}

\begin{tcolorbox}[breakable,colback=teal!4,colframe=teal!35!black,title={Before you move on}]
What does l'innamorato mean? (\textquotedblleft{}Sweetheart; in love\textquotedblright{}.) Say it once more.
\end{tcolorbox}
\section[la tenerezza]{la tenerezza — tenderness}

\label{lesson:IT-C491-tenerezza}



\begin{quote}
Before the new one: say the Italian for to court, then the Italian for sweetheart.
\end{quote}

\subsection*{You'll want to know: la tenerezza}
\textbf{la tenerezza} — \textquotedblleft{}tenderness\textquotedblright{}.

\begin{quote}
Mi guarda con tenerezza.
\end{quote}

\begin{quote}
\emph{She looks at me tenderly.}
\end{quote}

Built on \emph{tenero} \textquotedblleft{}tender, soft\textquotedblright{}.

\begin{cousinweb}
\emph{tenerezza} is built on \emph{tenero} \textquotedblleft{}tender\textquotedblright{}, inherited from Latin \emph{tener} \textquotedblleft{}soft, delicate\textquotedblright{} (Lewis \& Short), with \emph{-ezza} for a quality.
\end{cousinweb}

\begin{grammarlens}[title={What you've built}]
One more word for telling, explaining and describing at B1.
\end{grammarlens}

\subsection*{Your turn}
\begin{itemize}
\raggedright
  \item \emph{Say it:} \emph{la tenerezza}
  \item \emph{Say it:} \emph{la tenerezza}, once more
  \item \emph{Say it:} a sentence of your own with \emph{la tenerezza}
  \item \emph{Recall:} say the Italian for to court, then the Italian for sweetheart, then say \emph{la tenerezza} again
\end{itemize}

\begin{tcolorbox}[breakable,colback=teal!4,colframe=teal!35!black,title={Before you move on}]
What does la tenerezza mean? (\textquotedblleft{}Tenderness\textquotedblright{}.) Say it once more.
\end{tcolorbox}
\section[il bacio]{il bacio — kiss}

\label{lesson:IT-C491-bacio}



\begin{quote}
Before the new one: say the Italian for sweetheart, then the Italian for tenderness.
\end{quote}

\subsection*{You'll want to know: il bacio}
\textbf{il bacio} — \textquotedblleft{}kiss\textquotedblright{}.

\begin{quote}
Mi ha dato un bacio sulla guancia.
\end{quote}

\begin{quote}
\emph{She gave me a kiss on the cheek.}
\end{quote}

The plural is \emph{i baci}. At the end of a message: \emph{baci} — love.

\begin{cousinweb}
\emph{bacio} is inherited from Latin \emph{bāsium} \textquotedblleft{}a kiss\textquotedblright{} (EtymDB; Lewis \& Short).
\end{cousinweb}

\begin{grammarlens}[title={What you've built}]
One more word for telling, explaining and describing at B1.
\end{grammarlens}

\subsection*{Your turn}
\begin{itemize}
\raggedright
  \item \emph{Say it:} \emph{il bacio}
  \item \emph{Say it:} \emph{il bacio}, once more
  \item \emph{Say it:} a sentence of your own with \emph{il bacio}
  \item \emph{Recall:} say the Italian for sweetheart, then the Italian for tenderness, then say \emph{il bacio} again
\end{itemize}

\begin{tcolorbox}[breakable,colback=teal!4,colframe=teal!35!black,title={Before you move on}]
What does il bacio mean? (\textquotedblleft{}Kiss\textquotedblright{}.) Say it once more.
\end{tcolorbox}
